%% sections/10_discussion.tex

\section{논의}
\label{sec:discussion}

\subsection{llm-d와의 비교}

llm-d~\cite{llmd2025}는 Kubernetes 기반 LLM disaggregation으로
워크로드 인지 라우팅을 지원하는 강력한 접근이다.
\S\ref{subsec:res-b200}에서 동일 환경(B200$\times$8, 8개 모델)으로 실측한 결과,
llm-d 는 vanilla 대비 80\%(45/56 셀)에서 우위를 보이나
최선 단일 방법(suffix)을 능가하는 경우는 18\%(10/56)에 그쳤다.
즉 클러스터 라우팅의 추가 이득은 서버 내 spec-decode 최적 선택 대비 제한적이며,
이는 경량 워크로드 게이팅이 대부분의 가치를 포착할 수 있다는 본 연구의 전제를 뒷받침한다.
두 시스템의 설계 철학은 다음과 같이 대비된다.

\begin{table}[t]
\centering
\caption{\algname\ vs.\ llm-d 설계 비교.}
\label{tbl:vs-llmd}
\footnotesize
\begin{tabular}{@{}lll@{}}
\toprule
\textbf{항목} & \textbf{llm-d} & \textbf{\algname} \\
\midrule
배포 범위 & 클러스터(다중 노드) & 단일 서버 \\
인프라 요구 & Kubernetes, etcd, & vLLM $+$ FastAPI \\
            & 노드 간 KV 전송 & (코어 외부 라우팅) \\
라우팅 기준 & prefill/decode 분리 & 워크로드$\times$모델 \\
투기적 디코딩 & 별도 설계 필요 & spec-decode engine 내장 \\
CPU 활용 & 미다룸 & NUMA pin $+$ branchy \\
실측 이득 & vanilla 대비 80\% & oracle 게이트 상한 \\
(B200, 56셀) & 최선단일 대비 18\% & (모델$\times$조건 최적) \\
\midrule
목표 & 클러스터 규모 확장 & 서버 내 효율 극대화 \\
\bottomrule
\end{tabular}
\end{table}

두 접근은 상호 보완적이다.
대규모 클러스터에서 llm-d가 노드 간 prefill/decode를 분리하는 동시에,
각 노드 내부에서 \algname 이 워크로드 레벨 게이팅과 CPU 슬랙 하베스팅을 수행할 수 있다.

\subsection{Code Agent에 대한 함의}

\S\ref{subsec:mot-code-fail}에서 논의한 바와 같이,
code agent의 반복 호출 패턴은 투기적 디코딩 회귀의 누적 효과를 심화시킨다.
\algname 의 워크로드 분류기는 현재 단발성 코드 생성(code workload)과
code agent를 동일하게 처리한다.

% TODO: Code agent 특화 게이팅 전략.
%       (1) tool calling JSON 생성(반복 구조 → suffix 적합)과
%       (2) reasoning chain(고유 토큰 → suffix 불리)을 구분하는
%       세분화 분류기 필요.
%       현재 regex 분류기를 code agent 세부 패턴으로 확장하는 실험 필요.

\subsection{한계}

\textbf{하드웨어 특화.}
대규모 실 트레이스 평가는 DGX B200$\times$8 + Xeon 8570 에서,
CPU lever microstudy 는 개발 플랫폼 H100$\times$8 + Xeon SPR 에서 수행되었다.
처리량 패턴은 두 플랫폼에서 일관되나, CPU 하베스팅
hyperparameter($T_{\text{aux}}$, branchy 패턴)의 B200 재측정 및
다른 GPU 세대(A100, H200), AMD GPU, 단일 소켓 서버로의 일반화는 미완이다.
\algname 의 $T_{\text{gpu}}$ 자동 측정이 부분적 일반화를 지원하지만
multi-environment 검증은 후속 연구이다.

\textbf{Branchy 보조 작업의 미검증.}
branchy 보조 작업의 $+5.28\%$(1-run)는 multi-run binding 검증이 완료되지 않았고,
그 메커니즘(CPU active-state 유지, \S\ref{sec:mechanism})도 직접 측정으로
확증되지 않은 \emph{예비 관찰}이다.
따라서 \algname 의 CPU 하베스팅에서 production-ready로 권장하는 lever는
검증·설명이 명확한 NUMA 바인딩($+3.13\%$ warm-only)이며,
branchy 보조 작업은 향후 검증 대상이다.

\textbf{게이팅 순효과의 공정 측정.}
\S\ref{subsec:gating-origin}에서 분석했듯 현재 end-to-end 결과는
baseline-only(GPU 4장)와 \algname-gated(GPU 8장)의 GPU 수 차이로 인해
병렬화 효과를 포함한다.
게이팅 decision의 순효과를 공정하게 측정하려면
양 구성의 GPU 수를 동일하게 고정한 통제 실험이 필요하다(후속 연구).

\textbf{짧은 디코딩 길이의 CPU lever 측정.}
처리량 평가는 실제 서빙 trace(ShareGPT, WildChat, LMSYS-Chat, SWE-bench,
MBPP, HumanEval) 2{,}159 prompt 와 \texttt{max\_tokens}=8{,}192 로
대규모 재현되었으나(\S\ref{subsec:res-b200}),
CPU lever microstudy 는 아직 \texttt{max\_tokens}=32 의 짧은 출력 환경에서
수행되었다. 긴 디코딩(수천 토큰)에서의 lever 효과 및 실 trace 기반
CPU 하베스팅 재측정은 후속 과제이다.

\textbf{Gating 규칙의 확장성.}
현재 벤더-워크로드 게이팅 규칙(Qwen-72B code $\to$ baseline)은
수동 측정 결과에서 도출된다.
Llama-405B FP8, 새 아키텍처 등 미측정 모델에서
runtime accept rate 모니터링 기반 자동 규칙 생성이 필요하다.

\subsection{재현성}

실험 스크립트, raw benchmark JSON, \algname\ 참조 구현은
보조 자료\footnote{Anonymous; URL will be provided in camera-ready.}로 공개한다.
모든 참조의 publication metadata는 DBLP, OpenReview, arXiv, Semantic Scholar의
다중 채널 검증(2026-05-27 기준)으로 확인되었다.
